\section{Finite tidal observations and the category conclusion}
\label{sec:packet}

The comparison development constructed above supplies a branch along
which a bounded-curvature geodesic experiences an exponentially large
optical frequency. We now use that frequency amplification to make a
finite tidal observation arbitrarily large by an arbitrarily small
change in any prescribed finite list of data seminorms. The construction
has three distinct objects: an approximate interior metric perturbation
$h^p$, an approximate exterior perturbation $h^B$ transferring its entry
jets to the original initial end, and exact vacuum data $d_T$ obtained
by correcting the constraints of the exterior seed. A separate Cauchy
attachment places the observation in the full development of $d_T$.

All wave comparisons
in this section use ordinary wave-map reduction with the fixed vacuum
background $g=g_d$ as target. They are local hyperbolic comparisons,
independent of the exterior modulation gauge.
Explicitly, for the unknown metric $\widetilde g$ we use
\begin{equation}
 \begin{split}
 \operatorname{Ric}(\widetilde g)_{\mu\nu}
      +\nabla^{\widetilde g}_{(\mu}
          \Upsilon_{\nu)}(\widetilde g;g)&=0,\\
 \Upsilon_\mu(\widetilde g;g)
   &=-\widetilde g_{\mu\lambda}\widetilde g^{\alpha\beta}
       \bigl(\Gamma(\widetilde g)^\lambda_{\alpha\beta}
                           -\Gamma(g)^\lambda_{\alpha\beta}\bigr).
 \end{split}
 \label{pk2:local-reduction}
\end{equation}
The target $g$ is a fixed smooth coefficient field in this equation.
There are no spatially nonlocal gauge additions or modulation feedback
variables in \eqref{pk2:local-reduction}. Its exact wave-map Cauchy
initialization and the vacuum constraints give homogeneous gauge
propagation and ordinary finite domain of dependence.

\begin{proposition}[A finite tidal observation]
\label{pk2:finite-signal}
Let $d\in\mathcal U_\epsilon$, and let $w\in T\Sigma$ be a seed whose future
unit timelike geodesic has a noncorner branch as its first runout from
the comparison development $\mathcal P_d$. Suppose its curvature is
bounded in the parallel positive metric up to that runout. For every
integer $m\ge10$, every $\eta>0$, and every $J>0$, there are smooth
complete vacuum data $d_T\in\mathcal D$ and seeds $w_T\to w$, for
arbitrarily large finite $T$, such that
\[
 p_m(d_T-d)\longrightarrow0,
\]
and the unit geodesic of $w_T$ in the full MGHD of $d_T$ has a finite
point with tangent $U_T$ and a unit spacelike vector $X_T\perp U_T$
for which
\[
 \bigl|\operatorname{Riem}(g_{d_T})(U_T,X_T,U_T,X_T)\bigr|>J.
\]
In particular $p_m(d_T-d)<\eta$ for sufficiently large $T$, and the
data may be required to remain in any prescribed relatively open
neighborhood of $d$.
\end{proposition}

We prove the proposition in the following subsections. The orders of
all comparison norms are fixed before choosing the positive frequency
exponent; formal accuracy is chosen afterwards, and $T$ is chosen last.
The notation $e^{o(T)}$ denotes a bound by $C_\sigma e^{\sigma T}$ for
every $\sigma>0$, at each fixed derivative and formal expansion order.
Constants and thresholds may depend on those orders and on $d$.
For a positive quantity, ``two-sided $e^{o(T)}$'' means that the same
convention applies to its reciprocal. No uniformity in an unbounded
sequence of derivative orders is asserted.

\subsection{The branch state and the normalized coefficients}
\label{pk2:branch}

Arrange the double-null labels so that
\[
 u=y_0,\qquad v=y_1,\qquad
 s=\frac{u+v}{2},\qquad x=\frac{v-u}{2},
 \qquad v\longrightarrow\infty,\quad u\longrightarrow u_*<\infty
\]
on the chosen branch. The label-interchange construction proved above
provides these conventions also for the other branch, on the same
surface $\mathcal S=\{s=0\}$ and the same product wedge. Its new angular
labels are obtained by finite sphere flow at every finite point; their
estimates are then rederived from the prepared entry data. Thus no
uniform derivative bound for an unestimated long angular flow is used.
Throughout the experiment these are fixed background coordinates.

Write
\[
 g=-2a\,du\,dv+
 \gamma_{AB}(d\theta^A-b^Adu)(d\theta^B-b^Bdu),
 \qquad L_0=\partial_u+b^A\partial_A,
 \quad L_1=\partial_v,
 \quad a=q_0e^{-\kappa(u+v)}A.
\]
Let $P_T$ be the background geodesic point with $v=T$, and let $U$ be
its unit tangent. Set
\[
 \Xi=\dot\theta-b\dot u,\qquad p=a\dot v>0,
 \qquad P=\gamma\Xi,\qquad \pi_u=-p-b\cdot P.
\]
The dot here denotes proper-time differentiation; $P$ is the angular
covector momentum, whereas $P_T$ is a spacetime point.

\begin{lemma}[Branch normalization]
\label{pk2:branch-normalization}
Along the branch, $p\asymp1$, and $P$, $\Xi$, and $\dot u$ remain
bounded. In particular, with $Y_T=\dot v(P_T)$,
\begin{equation}
 Y_T\asymp e^{\kappa T}.
 \label{pk2:blueshift}
\end{equation}
\end{lemma}
\begin{proof}
Unit normalization gives $1+|\Xi|_\gamma^2=2p\dot u$. The geodesic
equation used in the corner-ratio argument implies
\[
 |\dot{\log p}|\le C(\dot u+|\Xi|_\gamma).
\]
The integral of $\dot u$ is finite. Also, by the same normalization
and Cauchy--Schwarz,
\[
 \int|\Xi|_\gamma\,d\tau
 \le C\left(\int du\right)^{1/2}
          \left(\int a\,dv\right)^{1/2}<\infty.
\]
Thus $p\asymp1$. Near the angular limit the unit-shell Hamiltonian
with $v$ as time is
\begin{equation}
 -\pi_v=\frac{a(1+|P|_{\gamma^{-1}}^2)}{2p},
 \qquad
 \frac{du}{dv}=\frac{a(1+|P|_{\gamma^{-1}}^2)}{2p^2}.
 \label{pk2:branch-hamiltonian}
\end{equation}
The low coefficient bounds give
$|dP/dv|\le C(p+|P|)\,du/dv$. Since the total variation of $u$ is
finite, Gronwall bounds $P$ and hence $\Xi$ and $\dot u$.
Finally $a\asymp e^{-\kappa T}$ at $P_T$, proving
\eqref{pk2:blueshift}.
\end{proof}

Fix $u_2>u_*$, $c_2>0$, and $c_0>0$ small relative to these margins.
The finite comparison domain is
\begin{equation}
 0\le s\le\frac{T+u_*}{2}+c_0,
 \qquad u+\frac{c_0s}{T}\le u_2,
 \qquad v+\frac{c_0s}{T}\le T+c_2.
 \tag{P1}\label{pk2:domain}
\end{equation}
Its slice lengths in $x$ have a positive lower bound. All constructions
are first made with enlarged domains and margins. For auxiliary
transports fix $u_f>u_2$. The strip $-v\le u\le u_f$, with $v$ in a
fixed neighborhood of $T$, is contained in the expanding region of
the higher double-null estimate, with bounded-direction endpoint
$u_f$. This observation includes the moving initial foot $u=-v$.

Use normalized vector slots and derivative fields
\[
 \hat e_i=a^{-1/2}L_i,\qquad \hat e_A=\partial_A,
 \qquad Z_\alpha=\sqrt a\,\hat e_\alpha,
 \qquad W_t=L_0+L_1,\quad W_x=L_1-L_0.
\]
Thus $Z_i=L_i$. For a covariant two-tensor entry $n$, let $n_i$ count
its radial $i$-slots and put $c_n=\kappa(n_0-n_1)/2$.
In the coefficient estimates below, $Db$ denotes the collection of
first ordinary coordinate derivatives of $b$ in the fixed
$(u,v,\theta)$ charts.

\begin{lemma}[Coefficient inputs for the finite comparison]
\label{pk2:coefficients}
The matrices $V_\alpha$ of $\nabla_{Z_\alpha}$ on the normalized
vector basis are bounded. In the bulk limit
$\min(v,-u,s)\to\infty$,
\[
 V_0\longrightarrow
 D_0=\operatorname{diag}(-\kappa/2,\kappa/2,0,0),
 \qquad V_1\longrightarrow-D_0,
 \qquad V_A\longrightarrow0.
\]
The contracted differentiated connection matrices in $a\square_g$
and every normalized entry of $a\operatorname{Riem}(g)$ are bounded
and small in that limit. The undifferentiated connection products in
the tensor wave operator retain the radial diagonal constants.
Moreover
\[
 b,\ Db,\ W_t\gamma,\ W_t\log a+2\kappa\longrightarrow0
\]
there, while first ordinary derivatives of $\log a$ and $\gamma$ are
bounded. Every fixed ordinary derivative of the normalized arrays
has size $e^{o(T)}$ on the required domains. Derivatives of $a$,
$\sqrt a$, and $\partial_vb$ retain, respectively, the factors
$a$, $\sqrt a$, and $a$.
\end{lemma}
\begin{proof}
Write $B_0=b$, $B_1=0$, and $j=1-i$.
The complete connection formulas include
\begin{align}
 \nabla_{L_i}\hat e_i&=\tfrac12\ell_i\hat e_i,&
 \nabla_{L_i}\hat e_j&=-\tfrac12\ell_i\hat e_j
                         +\sqrt a\,\xi_i^A\partial_A,
 \label{pk2:radial-connection}\\
 \nabla_{L_i}\partial_A
   &=(\chi_{i,A}{}^B-\partial_AB_i^B)\partial_B
                         +\sqrt a\,\xi_{i,A}\hat e_i,
 \notag\\
 \nabla_{Z_A}\hat e_i
   &=\chi_{i,A}{}^B\partial_B
      +\tfrac12\sqrt a(\xi_i-\xi_j)_A\hat e_i,&
 \nabla_{Z_A}\partial_B
   &=\sqrt a\,\Gamma_{AB}^C\partial_C
           +\chi_{0,AB}\hat e_1+\chi_{1,AB}\hat e_0.
 \notag
\end{align}
The low estimates and pure-jet limits give the asserted limiting
matrices. In the mixed trace the radial derivatives are cross
derivatives; in particular
$L_1\partial_Ab=\partial_A\partial_vb$ retains $a$.
For curvature use the exact matrix identity
\[
 aR(\hat e_\alpha,\hat e_\beta)
 =Z_\alpha V_\beta-Z_\beta V_\alpha
      +[V_\alpha,V_\beta]-V_{[Z_\alpha,Z_\beta]}.
\]
The brackets satisfy
\[
 [L_0,L_1]=\sqrt a(\xi_0-\xi_1)^AZ_A,
 \qquad
 [L_i,Z_A]=\tfrac12\ell_iZ_A-(\partial_AB_i^B)Z_B.
\]
The new unsuppressed pure longitudinal curvature term is
$L_i\chi_i$, covered by the pure-jet estimate. Longitudinal torsion
derivatives reduce to $\beta_i$, $d\ell_i$, and $\tau_i\xi_i$ by
the double-null identities. The actual low bounds on
$\ell_i+\kappa$ and $\chi_i$ in angular $H^4$ give $d\ell_i$ and
$\beta_i$ in $H^3$; the pure jets are in $H^2$, which embeds in
$C^0$ on the spheres. Thus the required limits do not differentiate
an uncontrolled moving profile. The factor $\sqrt a$ in the angular
connection and the limiting commutation relation do not create a
constant curvature term. The only constant tensor-wave contribution
is
\[
 -(L_0+C)(L_1-C)-(L_1-C)(L_0+C),
 \qquad C=\operatorname{diag}(c_n),
\]
whose $n$th entry is the radial part of
$e^{2c_nx}a\square_g^{\mathrm{scal}}e^{-2c_nx}$.
The higher double-null bounds, on the strip including $u_f$, give
all remaining fixed derivative estimates with the stated factors.
\end{proof}

\subsection{A constrained oscillatory packet in the deep region}
\label{pk2:deep-packet}

Let $\mathcal P$ be the covariant tensor wave operator with curvature
potential obtained by trace reversal of the linearized wave-map reduced
vacuum equation at $g$, and let $\mathcal C=\operatorname{div}_g$ on
these tensors. Linearized Bianchi gives
\[
 \mathcal C\mathcal P=\mathcal P_1\mathcal C,
\]
where $\mathcal P_1$ is the one-form wave operator. Put
\[
 k_v=g^{-1}dv=-a^{-1}L_0,
 \qquad K_vh=h(k_v,\cdot),
 \qquad \mathcal T=2\nabla_{k_v}+\square_gv.
\]
For a sufficiently large finite $L$, solve the transport and gauge
hierarchy
\begin{equation}
 i\mathcal Tp_j+\mathcal Pp_{j-1}=0,
 \qquad iK_vp_j+\mathcal Cp_{j-1}=0,
 \qquad p_{-1}=0,
 \label{pk2:hierarchy}
\end{equation}
on $-v\le u\le u_f$, with $v$ near $T$.

Multiplication of the first equation by $-a/2$ gives connection
transport along $L_0$ plus area expansion, with a known source
supplied by $a\mathcal Pp_{j-1}$. At fixed $j$, the unknown normalized
entries of $p_j$ form triangular blocks: the block of weight $c_n$
has same-weight transport $L_0+c_n$, with bounded additional
coefficients small in the bulk, and couples only to higher-weight
blocks. The boundary values for each $p_j$ are assigned as follows:
\[
\begin{array}{c|c|l}
\text{entries}&c_n&\text{boundary value}\\ \hline
00&\kappa&\text{gauge equation at }u=-v\\
0A&\kappa/2&\text{gauge equation at }u=-v\\
01&0&\text{gauge equation at }u=-v\\
AB&0&\text{free angular data at }u=-v\\
1A&-\kappa/2&\text{zero at }u=u_f\\
11&-\kappa&\text{zero at }u=u_f
\end{array}
\]
In this table $A,B$ are angular slots. At $j=0$ the free angular
data are a compact smooth bump times an angular trace-free
polarization of unit norm at the foot of the $v=T$ ray through
$P_T$, multiplied by a constant unit phase to be fixed below.
They vanish for $j>0$. Solve the blocks in descending weight,
retaining the coupled angular transport within each block.
Nonnegative-weight blocks are integrated from $u=-v$, and
negative-weight blocks from $u_f$ backwards; the constant diagonal
part has the favorable sign in each chosen direction.

The gauge contraction uses precisely the entries with a $0$-slot,
so the terminal choices do not prescribe extra gauge data. To check
that the initial gauge equations persist, set
$q_j=iK_vp_j+\mathcal Cp_{j-1}$. Conjugated divergence compatibility
gives
\[
 i\mathcal T_1q_j+\mathcal P_1q_{j-1}=0,\qquad q_{-1}=0,
\]
where $\mathcal T_1$ is the corresponding one-form transport. Induction
and the zero initial gauge jets give $q_j=0$ throughout the strip.

All amplitudes and fixed ordinary derivatives in normalized entries
have size $e^{o(T)}$. The one-flank interior estimate applies to the
whole strip: choose its bounded coordinate to be $y_0=u\le u_f$
and its other coordinate $y_1=v\le CT$, with any fixed $C>1$.
This includes both $u=-v$ at entry and $u=u_f$ at the terminal edge.
The angular vector field $b$ and its first angular derivatives have
arbitrarily small average outside bounded end portions. Differentiated angular flow
and inverse flow therefore have subexponential fixed jets; higher
coefficient derivatives enter only lower induction orders. The
normalized sources $a\mathcal Pp_{j-1}$ and the prescribed gauge
components $\sqrt a\,\mathcal Cp_{j-1}$ obey the same bounds.
Locality of each transport, including the backwards integrations,
keeps support in the flow tube of the initial bump. The initial
patch can have fixed small size in uniform charts on $\mathcal S$,
although its center may vary with $T$. The leading tensor is
trace-free and annihilates $k_v$. Its screen polarization at $P_T$
has two-sided $e^{o(T)}$ norm, since it is parallel transported times
an area factor with small-average logarithmic derivative.

Choose $X\perp U$ with $dv(X)=0$ and $g(X,X)=1$, using an angular
unit vector and a bounded multiple of $L_0$. Choose its direction so
that the diagonal screen polarization has absolute value
$\sigma_T$, of two-sided $e^{o(T)}$ size. Fix the constant phase
to make that contraction fully real at $P_T$. Let $C_w$ bound the
background unit tidal contraction on this seed; it is finite by its
bounded parallel curvature and its fixed initial boost. Choose
$J_{\mathrm{amp}}>0$ so that $2J_{\mathrm{amp}}>J+C_w+2$, and set
\begin{equation}
 \lambda_p=e^{\zeta T},\qquad
 A_T=\frac{4J_{\mathrm{amp}}}{Y_T^2\lambda_p^2\sigma_T},\qquad
 h^p=A_T\operatorname{tr.rev.}_g
       \operatorname{Re}\left(e^{i\lambda_pv}
                    \sum_{j=0}^L\lambda_p^{-j}p_j\right).
 \label{pk2:packet-amplitude}
\end{equation}
Here $\zeta>0$ will be chosen after all base norm orders. The normalized
reduced-equation residual at $g+h^p$, multiplied by $a$, is
\begin{equation}
 e^{o(T)}\bigl(A_T\lambda_p^{-M_1}
                     +A_T^2\lambda_p^{P_*}\bigr)
 \label{pk2:deep-residual}
\end{equation}
at those orders. Increasing the finite hierarchy makes $M_1$ arbitrarily
large, whereas $P_*$ depends only on the fixed base orders. The linear
phase and transport errors cancel successively,
and subtracting the linearization leaves quadratic relative-field
products in scaled covariant derivatives and smooth inverse-array
coefficients, including algebraic $a\operatorname{Riem}(g)$ terms.
There is no inverse-lapse loss in these normalized residual arrays.

\subsection{The exact deep comparison and its observation}
\label{pk2:deep-comparison}

The next proposition specifies exactly how accurately the entry jets
must be transferred from the original initial end. We first define
its mixed error norms using the background fields. For a positive
mass parameter $\varepsilon_0$, a normalized entry $f$ of weight
$c_n$ has reference slice density
\begin{equation}\label{pk2:reference-density}
 \frac12\left(|W_tf|^2+|(W_x-2c_n)f|^2
        +2a|d_\theta f|_\gamma^2+\varepsilon_0^2|f|^2\right)
                   dx\,d\operatorname{vol}_\gamma.
\end{equation}
For a normalized tensor array $F$, let $E_j(F;s)$ be the square root
of the sum of these integrals over the $s$-slice of
\eqref{pk2:domain}, with $f=D^\alpha F_n$, over all entries $n$
and ordinary multiindices $|\alpha|\le j$. Angular derivatives and
localization use the fixed finite sphere atlas. These are the
through-$j$ mixed energies below; $\varepsilon_0$ will be fixed
sufficiently small after the desired growth-rate tolerance.

\begin{proposition}[Deep comparison]
\label{pk2:deep-proposition}
Fix a mixed Sobolev comparison order $k\ge10$. For every sufficiently
small fixed $\delta_2>0$, one may fix $\varepsilon_0>0$ and then
$\nu_0>0$ sufficiently small as follows. Suppose the exact
reduced-equation Cauchy jets on $s=0$ satisfy the vacuum constraints
and wave-map gauge, agree with $g$ for $v<(1-\nu_0)T$, and differ
from the jets of $g+h^p$ by at most
\begin{equation}
 e^{-7\kappa T/2}
 \quad\hbox{in }H^{k+1}\times H^k(dx\,d\theta).
 \label{pk2:entry-accuracy}
\end{equation}
Then these data have an exact vacuum development $g'$ on
\eqref{pk2:domain}, with normalized error
$F=g'-g-h^p$ bounded in the through-$k$ mixed energies by
\begin{equation}
 C e^{(\kappa+\delta_2)T}
 \left(e^{-7\kappa T/2}+A_T\lambda_p^{-M_1}
                         +A_T^2\lambda_p^{P_*}\right).
 \tag{P7}\label{pk2:deep-error}
\end{equation}
The conclusion holds when this bound has a fixed exponential margin
below $A_T$ and the required base packet norms are sufficiently
exponentially small. The same estimate, with fixed buffers, controls
normalized $C^2$ error.
\end{proposition}
\begin{proof}
Bootstrap the
error energies by $e^{-(2\kappa+c')T}$ for a small fixed $c'>0$.
The relative cone change is exponentially small. Therefore the
side inequalities in \eqref{pk2:domain} are strict outflow, since
$du/ds,dv/ds\ge-$exponentially small on future causal vectors.
Finite propagation keeps the exact change supported in
$v\ge(1-\nu_0)T-o(1)$.

For $f=D^jF_n$, where $D$ denotes ordinary derivatives through order
$k$, the scalar operator retaining the diagonal drift is
\[
 e^{2c_nx}a\square_{g'}^{\mathrm{scal}}(e^{-2c_nx}f).
\]
Lemma~\ref{pk2:coefficients} gives bounded top-order residual
coefficients small in the bulk, with $e^{o(T)}$ coefficients only
on lower commuted energies. The potentially unscaled angular
commutator has the form
\[
 (Db)\partial_v(\partial_AD^{j-1}F)
\]
when one derivative hits $b\partial_A\partial_v$. The angular
derivative is inside the ordinary order-$j$ field, so its coefficient
is the small $Db$. With two or more coefficient hits the field is at
a lower energy level. A first hit on $a\gamma^{-1}\partial_\theta^2$
retains at least $\sqrt a$ against the scaled angular derivative;
higher hits again enter lower levels. The factor $\partial_vb$
retains $a$, and angular partitions cost $O(|b|+\sqrt a)$ at top
order.

After retaining the exact principal coefficients, each perturbative
commutation or lower product contains an additional packet-small or
bootstrap-small factor at an error placement. The highest error
order $j+1$ can lose $O(a^{-1/2})\le Ce^{\kappa T/2}$; further
error factors through half the total order $j+2$ are controlled by
the mixed $k$-order bootstrap and slice Sobolev embedding. Taking
$\zeta$ sufficiently small leaves an exponential margin to absorb
these factors. Smooth higher background derivatives cost only
$e^{o(T)}$.

Apply massive scalar stress energy to $e^{-2c_nx}f$ with mass squared
$\varepsilon_0^2/(2a)$ and multiplier $e^{4c_nx}W_t$. Up to an
exponentially small relative metric change, its slice density is
\eqref{pk2:reference-density}. Thus the exact-principal stress energy
is uniformly equivalent to the reference mixed norm just defined.
The limiting nonsource bulk terms are
\[
 -4c_nW_tf(W_x-2c_n)f
       -2\kappa a|d_\theta f|_\gamma^2
       +\varepsilon_0^2fW_tf.
\]
Here $W_tx=0$, $W_xx=1$, $W_t\log a\to-2\kappa$, and
$W_t\gamma\to0$. The brackets of $W_t$ with $W_x,Z_A$ are small
in the scaled basis except for the limiting term $-\kappa Z_A$.
All exits have favorable flux. Since $|c_n|\le\kappa$, fixing
$\varepsilon_0$ small and then a bulk threshold gives square-root
energy rate $2\kappa+$arbitrarily small in the bulk, and a bounded
rate elsewhere. The part of the support not entirely in that bulk
has time length $O(1)+O(\nu_0T)$, whereas the whole interval has
length $T/2+O(1)$. More explicitly, for a fixed bulk threshold $B$
the bulk is $\min(v,-u,s)\ge B$. The support has $v>B$ for large
$T$, the region $s<B$ has bounded length, and $-u<B$ implies
$2s=u+v>(1-\nu_0)T-B-o(1)$. Its possible remaining time length
is at most $\nu_0T/2+O(B+1)$.
A common nonnegative majorant therefore integrates
to $(\kappa+\delta_2)T$ after slightly reserving its strict margin.

For clarity, if $E_j$ are the successive error norms, their
inequalities have the form
\[
 E_j'\le r_TE_j+b_T\sum_{\ell<j}E_\ell+S_j,
 \qquad \int b_T\le e^{o(T)}.
\]
Factoring out $\exp(\int r_T)$ and integrating successively in $j$
costs only a fixed polynomial in $\int b_T$. This is
\Cref{pk:triangular-comparison}, and is still $e^{o(T)}$ at every
fixed $k$. The initial mixed jets follow from the reduced equation
with the same source and subexponential costs: $s=0$ is uniformly
spacelike in its nondegenerate entry coordinates. This proves
\eqref{pk2:deep-error}.

Its strict exponential improvement closes the bootstrap. Ordinary
hyperbolic restarts on enlarged domains within the outflow margins
give continuation, since for each finite $T$ all coordinates are
nondegenerate. Homogeneous wave-constraint propagation gives vacuum.
Adding fixed Sobolev buffers gives the final $C^2$ assertion.
\end{proof}

\begin{lemma}[Tidal size and backward state]
\label{pk2:observation}
Under the strict parameter margins of
Proposition~\ref{pk2:deep-proposition}, the point $P_T$ has a future
unit tangent $U'$ and a unit spacelike $X'\perp U'$ for $g'$ with
\[
 |\operatorname{Riem}(g')(U',X',U',X')|>J.
\]
The $g'$-geodesic with state $(P_T,U')$, traced backwards to
$\mathcal S$, has state converging to the background geodesic's
transverse entry state and stays within the available side margins.
\end{lemma}
\begin{proof}
Put $H=U/Y_T$. In the normalized primal basis, $H$ has size
$O(\sqrt a)$ and $X$ has size $O(1)$, while $dv(H)=1$ and
$dv(X)=0$. The phase-phase leading term in
$\operatorname{Riem}(g')(H,X,H,X)$ has absolute value
\[
 \tfrac12 A_T\lambda_p^2\sigma_T.
\]
All other packet terms cost
$e^{o(T)}(A_T\lambda_p+A_T^2\lambda_p^{P_*})$; comparison terms
cost \eqref{pk2:deep-error} times $e^{o(T)}$, enlarging the fixed
$P_*$ if necessary. Indeed express the connection difference using
first $\nabla^g$ derivatives of $g'-g$, then use its covariant
derivative, connection products, and inverse-metric products for
curvature. Multiplication by $a$ allows both differentiated slots
to use $Z=\sqrt a\,\hat e$ in normalized entries, with only
subexponential conversion costs, including hits on the scaling.
Evaluation on two $H$ and two $X$ slots cancels this curvature
scaling. The single order-zero amplitude supplies the displayed
leading term; trace reversal does not alter it, and the ordinary
second-metric-derivative formula gives the factor $1/2$.

All remaining terms are $o(Y_T^{-2})$ under the choices made below.
A normalized near-identity frame isometry sends $U,X$ to a
$g'$-orthogonal future/unit and spacelike/unit pair $U',X'$, with
relative frame error $O(A_Te^{o(T)})$. Its curvature slot error is
still $o(Y_T^{-2})$ at the $H$ scale. Multiplying by $Y_T^2$ and
using the choice of $J_{\mathrm{amp}}$ proves the strict tidal bound.

For the backward state, use \eqref{pk2:branch-hamiltonian} on a
compact $u,\theta,\pi_u,P$ state tube around the background tail,
with $p=-\pi_u-b\cdot P$ bounded below. Its background Lipschitz
cost is $ae^{o(v)}$, integrable as $v\to\infty$. The perturbed unit
shell is solved on the same branch by putting $\pi_v=az'$: the
background equation for bounded $z'$ has derivative bounded away
from zero. Normalized covector arguments have size $O(a^{-1/2})$,
but the differentiated relative metric change is exponentially
smaller than $a$ on this bounded-$u$ path. The Hamiltonian and its
first state derivatives therefore change by exponentially small
absolute amounts. The terminal errors are
\[
 \delta(\pi_u,P)=O(A_Te^{o(T)}/\sqrt a),
 \qquad \delta(\pi_v/a)=O(A_Te^{o(T)}/a).
\]
The second estimate uses $\partial_v=\sqrt a\,\hat e_1$ and also
shows that the chosen future tangent lies on the continued root.
Differentiating the shell equation in state variables costs at most
the same inverse $a$ on normalized coefficient differences, with
subexponential background factors. Integrating the exponentially
small vector-field error over length $O(T)$ and applying Gronwall
with the integrable background Lipschitz majorant proves convergence
through the transverse hit of $s=0$. Fixed earlier portions use
ordinary compact geodesic flow. Uniform state and side margins keep
the perturbed paths in the stated domain.
\end{proof}

\subsection{Return rays to the original initial end}
\label{pk2:return-rays}

The interior packet is useful only if its entry jets can be produced
by data on the original bridge. We first verify the real-ray geometry
needed for the exterior beam construction, including its dependence
on the fixed spin.

Use the pertinent exterior time $t_e$, co-coordinate $y$, and lab
coordinate $z=y+C_e(t_e)$. Orient the aligned Boyer--Lindquist time
$T_K$ on this end so that the preparation parameter $I$ satisfies
\[
 I'=\lambda_-\varsigma>0,
 \qquad 0<c_I\le I'\le C_I,
 \qquad I(t_T)=T.
\]
Then $t_T\asymp T$. The launch patch near $v=x=T$ on $s=0$ has
$t_e=t_T+O(1)$. A fixed $C_2$ with margins covers its required
initial interval by $t_e\le t_T+C_2$. The part
$v<(1-\nu_0)T$ is earlier by at least $c\nu_0T$; the bounded joins
do not grow with $T$. The inner exterior exit lies strictly deeper
than this end of $\mathcal S$, with a further fixed background margin.

\begin{proposition}[Backward escape and flow bounds]
\label{pk2:escape}
The null rays launched by $dv$ on a sufficiently small fixed entry
patch, followed towards the past, return to the original initial
end. For every prescribed $\alpha>0$, after choosing a sufficiently
large fixed radius and then taking $T$ sufficiently large, all their
arrival points satisfy
\begin{equation}
 (1-\alpha)t_T<|z-z_*|<(1+\alpha)t_T,
 \qquad z_*=C_e(t_T).
 \tag{E1}\label{pk2:arrival-shell}
\end{equation}
Their entry crossing is unique and transverse within the computing
region, with separation from the entry level outside a uniform
collar. All fixed jets of the phase-space flow on all subsegments,
including launch-center derivatives, have bounds $e^{o(T)}$.
Tensor parallel transport has bounded norm and inverse along these
rays. The assertions hold with short past and future launch margins,
including passage through the deeper inner exit.
\end{proposition}
\begin{proof}
On compact spatial subsets of the co cylinder,
$g\to G_{p_\infty}$ in the required low differentiated orders.
Indeed $p\to p_\infty$ follows from the integrable modulation
derivative, and $p',c,c'$ and further required low jets tend to zero
by the positive-weight endpoint and feedback bounds. The compact
$A_H$-weighted endpoint estimates control $h_{\mathrm{co}}$, while
the target memories have bounded length there. Higher compact and
far estimates supply fixed jets of size $e^{o(T)}$ for time and
radius $O(T)$.

The entry surface is strictly spacelike and optically transverse.
A small uniform collar obtained by local eikonal and angular flow
therefore has bounded low transition jets and subexponential high
jets. The covectors $dv$ approach the frozen covectors
\[
 \lambda_-d(r_*+F+T_K)
\]
in aligned Kerr variables. This follows from the prepared-coordinate
Jacobian comparison, the decay of the parameter derivatives, and the
noncharacteristic null root. Accumulated phase offsets do not affect
this assertion. Here $F$ is the exact axisymmetric optical-profile
function, not the metric comparison error.

The frozen covector has zero angular Killing momentum. By Kerr
separation its radial potential divided by squared Killing energy is
\begin{equation}
 (r_K^2+\mathfrak a^2)^2
   -\Delta\left(|\partial_\vartheta F|^2
                         +\mathfrak a^2\sin^2\vartheta\right)_{\rm init}.
 \label{pk2:radial-potential}
\end{equation}
The bracket is at most $\mathfrak a^2+O(q_0^2)$, since
$F(r_-)=0$ and the profile is smooth. For $r_K\ge r_+$,
\[
 \Delta<r_K^2,
 \qquad
 \frac{(r_K^2+\mathfrak a^2)^2}{r_K^2}\ge4\mathfrak a^2.
\]
For every fixed $\mathfrak a>0$, choosing the original exact-profile
depth sufficiently small over a small physical parameter range
makes \eqref{pk2:radial-potential} strictly positive. It is also
positive between the horizons, where $\Delta<0$, with uniform
margins on compact crossings. This depth choice was made before the
exterior constants and packet parameters.

In the frozen chart $g^{-1}dv$ is past-directed with Killing energy
$-\lambda_-$. Its radial direction is increasing, since initially
\[
 \Sigma_K\,dr_K(g^{-1}dv)
   =\lambda_-\bigl(r_K^2+\mathfrak a^2+\Delta F_r\bigr)>0
\]
by the finite Hamilton--Jacobi root. At the outer horizon its radial
covector component has leading coefficient
$\lambda_-(r_K^2+\mathfrak a^2)/\Delta$ with bounded remainder.
The coordinate $T_K+r_*$ and the ingoing azimuth are regular, and
the angular momentum stays bounded, with zero azimuthal coefficient.
Thus the ray crosses the future event horizon transversely and
spends bounded regular $t_e$-time through every fixed compact
crossing. The same holds forwards from entry to the deeper exit.
No timelike stationary Killing field is needed in an ergoregion.

Near entry, $Q=(r_K-r_-)f$ with $f>0$ smooth and axisymmetric.
Separation and the strict radial potential bound give bounded
$|d\vartheta/dr_K|$. Hence $dQ/dr_K$ has a positive lower bound
in a sufficiently small profile neighborhood of $r_-$. Singular
azimuthal winding does not differentiate $Q$; round colatitude
length or smooth sphere charts give the same conclusion at poles.
Outside this neighborhood the used compact crossing portions are
separated from the entry level. Compact convergence of the metric
and prepared level graphs therefore gives a unique transverse entry
crossing, separated from all other used portions by a uniform collar.
Only times $t_T+O_R(1)$ are needed until a fixed large radius $R$
is reached. This comparison does not require the launching azimuthal
phase to converge.

The frozen branch is asymptotically radial. The limiting co-chart
is asymptotically affine in differential order, with small changes
of axes, shear, and boost. At a sufficiently large fixed co-radius
$R$, followed by sufficiently large $T$, the past-directed lab
velocity $\mathbf w=dz/d(-t_e)$ has length
$1+o_{R\to\infty}(1)$ and makes a small angle with $y/|y|$.
There $z=z_*+O_R(1)$ and $t_e=t_T+O_R(1)$, with uniformly
nondegenerate momentum and velocity normalizations on the whole
launch patch. Small drift and shear use only the fixed base
smallness, not frequency-dependent smallness.

In the far lab region, with emitter radius $r\asymp|y|$ and
$z=C_e(q)+rn$, $t_e=r+q$, the far estimates give, uniformly in time,
\begin{equation}
 \begin{split}
 |g-\eta|&\lesssim r^{-1+.01},\\
 |\partial g|&\lesssim r^{-2+.01}
   +r^{-1+e/2}\left\langle q/r^{\delta_f}\right\rangle^{-A_H/2},
 \qquad |\partial^2g|=o_{r\to\infty}(1).
 \end{split}
 \label{pk2:far-ray-coefficients}
\end{equation}
These are ordinary lab derivatives. To obtain them, use two spare
rotations in the point bounds and
$\partial_i=\ell_{n,i}T_l+r^{-1}b_iZ$. Second derivatives expand
with lab times commuted innermost; angular and $Z$ terms have an
extra $r^{-1}$, and the remaining low unrestricted jets are at worst
a bounded coefficient times the displayed bad-field size. The
differentiated emitter conversion coefficients remain bounded.
The smoothed target's differentiated symbol terms also tend to zero.

While $\mathbf w$ remains close to its crossing value,
$dy/d(-t_e)=\mathbf w+s_e(t_e)$ lies in an outward cone, $n$ is
close to $y/|y|$, and
\[
 \frac{dr}{d(-t_e)}
  =\frac{n\cdot(\mathbf w+s_e(q))}{1-n\cdot s_e(q)}
\]
has positive lower and upper bounds. Thus $q$ decreases at a rate
bounded below. On each large radial dyad $r\sim R_1$, integration
in $q$ over an interval of length $O(R_1)$ bounds the second term
of \eqref{pk2:far-ray-coefficients} by
\[
 O\left(R_1^{(e-(1-\delta_f)A_H)/2}\right).
\]
Here $A_H<2$, and the exponent is strictly negative for the fixed
far parameters. The first term is summable as well. The lab
geodesic equation in time parameter bounds both the velocity
derivative and logarithmic frequency-scale derivative by
$C|\partial g|$; spacelike temporal hyperbolicity bounds
$|\mathbf w|$. Choosing $R$ large makes the integrated cost small,
closes the outward-cone argument, and keeps momentum nondegenerate
until $t_e=0$. The ray then remains in the far portion on which
$t_e=0$ is the original bridge. Integrating its velocity proves
\eqref{pk2:arrival-shell}.

Finally consider the time-dependent null Hamiltonian root in lab
coordinates. In its first variation, the diagonal blocks and the
position-to-momentum block tend to zero with radius by the first
two derivative bounds in \eqref{pk2:far-ray-coefficients}; the
momentum-to-position block is bounded. On the rescaled variables
$(\epsilon_1\delta z,\delta\pi_{\rm sp})$, choose $\epsilon_1>0$
small first and then a sufficiently large fixed radius. The rate
is then arbitrarily small on the far portion. The complement has
time length $O_R(1)$. This gives subexponential all-subsegment
propagators. Higher jet induction has the same homogeneous
propagator and sources made from lower jets and the available
subexponential coefficient derivatives. It also covers derivatives
in the launch-center parameters. The integrated connection bound
gives bounded tensor parallel transport and inverse.
\end{proof}

\subsection{Constrained complex beams and matching}
\label{pk2:beams}

We now transfer the interior packet to the bridge using complex
Gaussian beams. Their complex phases provide transverse localization
through real caustics, while superpositions recover oscillatory Cauchy
data \cite{Ralston1982,Tanushev2008,LiuRunborgTanushev2013}.
We retain the phase-space and matching calculations because we also need
bounds on growing time intervals and Einstein wave-gauge compatibility.

\begin{proposition}[Exterior beam approximation]
\label{pk2:beam-proposition}
For every fixed base-order buffer $N_b$, every prescribed inverse
accuracy $M_1$, and sufficiently high finite phase and amplitude
orders, there is a real exterior perturbation $h^B$ such that
\begin{equation}
 \begin{split}
 \|h^B\|_{C^{N_b}}&\le A_T\lambda_p^{N_b+3/2}e^{o(T)},\\
 \text{reduced-wave and gauge residuals}
   &\le e^{o(T)}
       \left(A_T\lambda_p^{-M_1}+A_T^2\lambda_p^{P_0}\right),
 \qquad P_0=2N_b+7,
 \end{split}
 \tag{E2}\label{pk2:beam-estimates}
\end{equation}
after enlarging $N_b$ for the required trace and derivative buffers.
Its jets match those of $h^p$ at $\mathcal S$ with linear error
$A_T\lambda_p^{-M_1}e^{o(T)}$. The positive exponent $P_0$ depends
only on the base orders, not on later formal accuracy. On the
original initial slice its support lies in
\eqref{pk2:arrival-shell}, up to arbitrarily thin tube margins.
\end{proposition}
\begin{proof}
The construction has four stages. We first propagate complex phase
jets and the central tensor amplitude along the return rays. We then
impose the wave-gauge equations on every finite jet. At each inverse
frequency order, Gaussian superposition leaves a new central amplitude
subject to the gauge equation. Choosing its free part to match the
interior coefficient makes the real transport hierarchy determine the
matching normal jets.
Finally we choose finite Taylor orders and cutoffs to obtain the
required residual bounds.

\paragraph{Phase and central amplitude transport.}
Using the ray, subsegment, and collar bounds just established,
parameterize centers $Y$
by a fixed patch of the entry slice. With spatial entry coordinates
$w_s=(x,\theta)$, start the phase at
\[
 v(0,w_s)+\frac{i}{2}|w_s-Y|^2,
\]
choosing its temporal root to agree with $dv$ at the center.
Evolve the phase jets along the real rays in lab time by the
complex gradient-graph construction. Transverse projection from
the entry slice gives a positive imaginary Hessian on the lab
slice; its spacetime null direction is the ray direction.
For the position and momentum differential blocks $P_c,Q_c$,
the preserved symplectic identity is
\[
 P_c^*Q_c-Q_c^*P_c=2i\operatorname{Im}M_c,
\]
where $M_c$ is the initial complex Hessian. The all-subsegment
bounds imply $P_c^{-1}=e^{o(T)}$ in bound. The propagated Hessian
has positive imaginary part with inverse bounded by $e^{o(T)}$.
Taylor composition and reversion use only finite smooth jets at
real positions, so no spatial analyticity is required.

The central amplitude propagator is tensor parallel transport
times the inverse square root of the density ratio formed from
\[
 \sqrt{|\det g|}\,k^{t_e}\det P_c,
 \qquad k=g^{-1}d\varphi.
\]
It and its inverse have subexponential all-subsegment bounds.
Higher transport jets have the same homogeneous propagator with
lower-jet subexponential sources.

\paragraph{Compatibility of the wave and gauge jets.}
Impose both wave transports and gauge equations in the finite
Taylor quotient, taking the eikonal orders sufficiently high.
For a complex phase $\varphi$ set
\[
 \mathfrak h=g^{-1}(d\varphi,d\varphi),\quad
 \mathcal P_\lambda=\mathcal P+i\lambda\mathcal T-\lambda^2\mathfrak h,
 \quad\mathcal C_\lambda=\mathcal C+i\lambda K,
 \quad\mathcal P_{1,\lambda}=\mathcal P_1+i\lambda\mathcal T_1-\lambda^2\mathfrak h.
\]
Here $K$ contracts by $g^{-1}d\varphi$ and $\mathcal T,\mathcal T_1$
are the tensor and one-form transports. Conjugation gives the exact
identity $\mathcal C_\lambda\mathcal P_\lambda
=\mathcal P_{1,\lambda}\mathcal C_\lambda$. In the finite jet quotient
where all required eikonal jets vanish, put
\[
 W_j=i\mathcal T r_{j,Y}+\mathcal P r_{j-1,Y},\qquad
 q_j=iKr_{j,Y}+\mathcal C r_{j-1,Y},\qquad r_{-1,Y}=0.
\]
The coefficient of $\lambda^{2-j}$ is
\begin{equation}\label{pk2:beam-defect-identity}
 iKW_j+\mathcal C W_{j-1}=i\mathcal T_1q_j+\mathcal P_1q_{j-1},
 \qquad W_{-1}=q_{-1}=0.
\end{equation}
The extra phase degrees eliminate every differentiated $\mathfrak h$
term in these quotients. Once lower defects vanish, solving $W_j=0$
makes $q_j$ satisfy homogeneous one-form transport.
At launch the contraction by $k$ has the right inverse
\[
 f\longmapsto l\otimes f+f\otimes l-f(k)l\otimes l,
 \qquad l(k)=1.
\]
The nonzero pairing is inverted also in formal jets. This fixes
the forced gauge components while leaving free projected kernel
jets for matching.

\paragraph{Triangular matching at the entry slice.}
Integrate the beams with measure
$(\lambda_p/2\pi)^{3/2}dY$. In a short entry collar, center
positions project diffeomorphically, with the stated low and
subexponential high bounds. At a center, the phase difference
from $v$ starts quadratically and has positive imaginary Hessian.
Changing the center variable to its position on the current
entry slice and evaluating Gaussian moments gives an expansion
\[
 e^{i\lambda_pv}\sum_j\lambda_p^{-j}\mathfrak p_j
\]
in integer inverse powers. On the launch surface, set
$a_j(Y)=r_{j,Y}(0,Y)$. The initial phase and normalization give,
after removing $e^{i\lambda_pv}$,
\[
 \left(\frac{\lambda_p}{2\pi}\right)^{3/2}
 \int e^{-\lambda_p|w_s-Y|^2/2}r_{j,Y}(0,w_s)\,dY
 =a_j(w_s)+O(\lambda_p^{-1})e^{o(T)}.
\]
This holds in each required fixed finite spatial norm once
sufficient jets have been supplied; tube cutoffs contribute only
Gaussian-small errors. The integrand includes the chosen off-center
amplitude jets. At order $j$, the integrated coefficient is therefore
the new central value $a_j$ plus terms already fixed by earlier jets.
Derivatives of that new central amplitude first enter one inverse
power later, since odd Gaussian moments vanish.
For each fixed $T$, let an auxiliary frequency tend
to infinity in the differentiated wave and gauge residuals, after
supplying sufficiently many finite jets. Comparison of coefficients
shows that the integrated coefficients satisfy both real-phase
hierarchies; still unmatched higher jets can be completed subject
to their gauge conditions without changing a fixed coefficient.
Once lower coefficients match $p_i$, $i<j$, their transports match
in the collar, including their normal derivatives. The forced gauge
terms at order $j$ therefore agree, so the remaining discrepancy
lies in the contraction kernel on the initial slice. The free
central value cancels it. Extend that kernel value in spatial jets
by applying $\operatorname{Id}-R_kK$ to constant-component choices,
where $R_k$ is the displayed right inverse. These local projected
extensions preserve bump support and subexponential bounds.
The real-phase transport is noncharacteristic at the entry slice.
Once its sources and the new coefficient value agree, it fixes the
same normal derivatives on both sides. Thus matching the free central
value matches the whole collar and both required Cauchy jets; these
are not two independent conditions on that value.
This closes the constrained triangular induction.

\paragraph{Finite orders, localization, and residuals.}
All expansions are finite. For last amplitude index $L$, desired
spatial degree $D$, and a buffer $k_0$ for the required mixed-time jets,
it suffices to choose
\[
 N_j=D+2k_0+2(L-j)+2,\qquad 0\le j\le L,
\]
after enlarging $D,k_0$ for the finite Gaussian coefficient calculations.
The phase is supplied through at least $N_0+2$ and farther for the
required differentiated eikonal remainders. A wave source uses two jets
of the preceding amplitude, a gauge source one; this backwards choice
therefore supplies every needed source and matching jet. No bound
uniform in $L$ is required. Realize the lab Taylor beams with smooth cutoffs
in tubes of radius $e^{-b_pT}$, where $0<b_p<\zeta/2$. They may
be stopped smoothly after leaving the computing domain through
its strict inner exit. A residual Taylor power $d$ gains
$\lambda_p^{-d/2}e^{o(T)}$ under Gaussian localization. Tube-cutoff
errors are Gaussian-small, and each fixed differentiation loses
only a fixed positive power. The last transport remainder improves
with its finite hierarchy order. For differentiated integral
remainders in the collar, first restrict to
$|w_s-w_{\mathrm{center}}(s)|\lesssim\lambda_p^{-1/2+\epsilon_2}$,
with $\epsilon_2>0$ small for the specified finite order; the
complement is Gaussian-small. The Taylor remainder uses just
the supplied smooth subexponential jets.

There are no other entry crossings by
Proposition~\ref{pk2:escape}, and the far escape is disjoint from
the prepared surfaces. Each new undifferentiated formal jet costs
only $e^{o(T)}$, irrespective of the chosen finite degree.
Absolute center integration costs $\lambda_p^{3/2}$, giving the
first bound in \eqref{pk2:beam-estimates}. Nonlinear differences
cost at worst the square of that base bound with two extra
derivatives, giving $P_0=2N_b+7$ in
\eqref{pk2:beam-estimates}.
Increasing formal accuracy adds only finitely many Taylor and
transport coefficients, each with subexponential fixed jets and
its assigned nonpositive frequency power. At a fixed base
derivative order it does not increase the differentiation loss
$\lambda_p^{N_b}$ or the integration loss $\lambda_p^{3/2}$.
The same observation in the interior hierarchy fixes $P_*$ from
its comparison orders before $M_1$ is increased. Constants and
the eventual threshold for $T$ may increase with formal accuracy.
Take real parts, undo trace reversal, and multiply by $A_T$ to
obtain $h^B$. Enlarged domains through $t_T+O(1)$ and radius
$C(1+T)$ contain all tubes and their buffers. At $t_e=0$ the tubes
lie in the shell \eqref{pk2:arrival-shell}, where the slice is the
original bridge. Gauss--Codazzi and the differentiated wave and
gauge residuals show that its induced data, set equal to $d$
elsewhere, have a compactly supported constraint defect with
the residual bound \eqref{pk2:beam-estimates}.
\end{proof}

\subsection{Exact constraints outside a protected ball}
\label{pk2:constraints}

The beam seed is only approximately constrained. The following
precise correction result permits exact vacuum data without fixing
the asymptotic charges or disturbing a large interior region.

For a tensor in Cartesian components define
\[
 \|f\|_{H_i^j}
 =\|f\|_{H^j(B_2)}
   +\sup_{L\ge1}L^i
       \|f(L\,\cdot)\|_{H^j(1<|\cdot|<2)},
 \qquad
 X^s=H_1^{s+2}\oplus H_2^{s+1}.
\]
Let $\Phi$ denote the vacuum scalar and momentum constraint map.

\begin{proposition}[Protected exterior correction]
\label{pk2:correction-theorem}
Fix an integer $s\ge3$. Suppose $d_B=(\delta+a_B,b_B)$ are smooth data on
$\mathbb R^3$, the metric is positive, and all derivatives of
$a_B,b_B$ have symbol orders $-1,-2$, respectively. Suppose a
smooth compact source $\phi$, supported in $\{|y'|\ge1\}$,
agrees with $\Phi(d_B)$ there. If
\[
 \|a_B\|_{H_1^{s+2}}+\|b_B\|_{H_2^{s+1}}
       +\|\phi\|_{H_4^s}
\]
is sufficiently small, there is a smooth correction $v_c$ supported
in $\{|y'|\ge1\}$ such that
\[
 \Phi(d_B+v_c)=0\quad\text{there},
 \qquad \|v_c\|_{X^s}\le C_s\|\phi\|_{H_4^s}.
\]
The correction is smooth across the support boundary, preserves
positivity, and satisfies symbol bounds at every order. There is
no higher-order smallness assumption and no condition on the
artificial interior defect. No moment condition is imposed on
$\phi$; the correction may have a noncompact tail changing the
asymptotic coefficients.
\end{proposition}
\begin{proof}
Apply \Cref{pk:constraint-nonlinear} with $Y^s=H_4^s$ and
$X^s=H_1^{s+2}\oplus H_2^{s+1}$. Its right inverse fixes the
unit inner ball and permits the stated asymptotic tail. Its single-base-order
nonlinear solve proves both the estimate and higher smoothness without
additional smallness assumptions.
\end{proof}


Apply this proposition with
\begin{equation}
 \rho_T=(1-2\alpha)t_T,
 \qquad y'=\frac{z-z_*}{\rho_T}.
 \label{pk2:rescaling}
\end{equation}
Write the unperturbed datum on this end as $d=(h,K)$.
This rescaling keeps the Cartesian metric component functions and
multiplies the second-form component functions by $\rho_T$:
\begin{equation}\label{pk:constraint-rescaling}
 \widetilde h(y')=h(z_*+\rho_Ty'),\qquad
 \widetilde K(y')=\rho_TK(z_*+\rho_Ty').
\end{equation}
Both constraint components then scale homogeneously by $\rho_T^2$.
Small fixed
drift and small $\alpha$ give $|z|\asymp\rho_T|y'|$ for
$|y'|\ge1/2$. More explicitly, the original symbol bounds imply,
for every fixed derivative order $j$,
\begin{equation}\label{pk:constraint-background}
 |\partial_{y'}^j(\widetilde h-\delta)|
       \le C_j\rho_T^{-1}|y'|^{-1-j},\qquad
 |\partial_{y'}^j\widetilde K|
       \le C_j\rho_T^{-1}|y'|^{-2-j}.
\end{equation}
The unperturbed scaled pair therefore has $X^s$ size
$O_s(\rho_T^{-1})$. Cut these perturbations off smoothly to flat
data between $|y'|=3/4$ and $1/2$; no constraint condition is
needed on this artificial interior.

By \eqref{pk2:arrival-shell}, the scaled beam seed lies strictly
outside the unit ball. Its size at the chosen solve order tends
to zero after $\zeta$ is made sufficiently small, with only
polynomial rescaling costs. The exterior constraint defect,
extended by zero, is smooth and compactly supported because
the seed vanishes near the unit sphere and the remaining data
there are exactly vacuum. Its norm has the residual bound in
\eqref{pk2:beam-estimates}, with polynomial factors included in
$e^{o(T)}$. Thus every hypothesis of
Proposition~\ref{pk2:correction-theorem} holds for large $T$.

Undo the scaling and retain $d$ on the protected interior and
the rest of the bridge. The correction is smooth and identically
zero on the protected side. The resulting data $d_T$ are smooth,
vacuum, positive, and complete; outside a compact set their metric
is uniformly comparable to the original complete metric. They
belong to $\mathcal D$ at every order, agree with $d$ on
$|z-z_*|<\rho_T$ and off the selected end, and may change the
permitted asymptotic coefficients. Here the Cartesian symbol norms
on the end are equivalent, at every fixed order, to the
fixed-background covariant norms defining $\mathcal D$.

For the prescribed $m\ge10$, choose $s\ge m+3$, increasing it also
for the comparison and trace buffers. Scaled annular Sobolev
embedding and unscaling give
\begin{equation}
 p_m(d_T-d)
 \le e^{o(T)}\left(
       A_T\lambda_p^{m+5/2}
       +A_T\lambda_p^{-M_1}+A_T^2\lambda_p^{P_0}\right).
 \label{pk2:data-approximation}
\end{equation}
One extra beam derivative accounts for the second form. To see the
unscaling cost for the correction, write its scaled components as
$(q_c,p_c)$. Its physical components are $q_c(y')$ and
$\rho_T^{-1}p_c(y')$. Thus a derivative of order $j$, multiplied
by the respective physical weights $|z|^{1+j}$ and $|z|^{2+j}$,
costs at most $C\rho_T$ times the corresponding scaled weighted
bound in both cases. These polynomial factors are harmless.
No compact-support assertion
about the final exact change at infinity is used.

\subsection{Ordinary exterior comparison and causal protection}
\label{pk2:exterior-comparison}

The corrected data are now exact vacuum data. We next show that their
ordinary wave-map evolution supplies the entry jets in
\eqref{pk2:entry-accuracy}. Give $d_T$ to the partial exterior problem
on $t_e=0$. Its bent compact portion lies in the dependence region
of unchanged data for large $T$. Choose the initial lapse and shift
as for $g+h^B$, and determine the normal metric derivatives from
the actual corrected second form and the exact wave-map condition.
Spacelikeness makes this algebraic system invertible. The resulting
metric and first jets differ from $g+h^B$ by the residual bound
in \eqref{pk2:beam-estimates}, by the correction and gauge-defect
estimates, and agree exactly with $g$ in the protected region.

Take the exterior comparison through $t_e=t_T+C_2$, with the fixed
deeper inner exit and shrinking outer cutoff
\begin{equation}
 |y|<R_T(t_e),\qquad
 R_T(t_e)=C_{\mathrm{out}}(1+T)
              +V_{\mathrm{out}}(t_T+C_2-t_e),
 \label{pk2:outer-cutoff}
\end{equation}
where $C_{\mathrm{out}}$ includes all needed spatial buffers and
$V_{\mathrm{out}}$ exceeds the radial causal-speed bound.

\begin{proposition}[Exterior transfer with exact early agreement]
\label{pk2:exterior-proposition}
For any prescribed small $\delta_3>0$, the localization parameters
and far threshold can be chosen so that the exact ordinary
wave-map solution exists on this finite exterior domain and its
error from $g+h^B$, in the required fixed mixed norms, is bounded by
\begin{equation}
 C e^{\delta_3T}
       \left(A_T\lambda_p^{-M_1}+A_T^2\lambda_p^{P_0}\right).
 \tag{E3}\label{pk2:exterior-error}
\end{equation}
It is exactly $g$ on the earlier needed compact-in-radius region
up to $t_T-\alpha'T$, where $\alpha'>0$ can be prescribed
sufficiently small. In particular, with the choices below, its
entry jets agree exactly with $g$ for $v<(1-\nu_0)T$ and otherwise
meet \eqref{pk2:entry-accuracy}.
\end{proposition}
\begin{proof}
We first prove the exact support assertion; the correction's tail
need not be compact. Write $s_* =\sup|C_e'|<1$, fixed by the base
smallness. Suppose a causal path from the changed initial support
first reaches a fixed sufficiently large co-radius $R$ at time
$\tau\le t_T$. Before this hit its lab speed is at most
$1+\delta_R$, where $\delta_R\to0$ as $R\to\infty$, also in the
exponentially small comparison bootstrap. Its initial point is
outside the radius $\rho_T$ about $z_*$, and its hit point is
$C_e(\tau)+O(R)$. Thus
\begin{equation}
 (1-2\alpha)t_T
 \le(1+\delta_R)\tau+s_*(t_T-\tau)+O(R).
 \label{pk2:first-hit}
\end{equation}
Equivalently,
\begin{equation}
 t_T-\tau\le
   \frac{2\alpha+\delta_R}{1+\delta_R-s_*}\,t_T+O_R(1).
 \label{pk2:late-first-hit}
\end{equation}
Since $t_T\asymp T$, a fixed large $R$ and a fixed small $\alpha$
make the right side less than any prescribed small multiple of
$T$ for large $T$. Only the inner exclusion radius of the change
enters this estimate; its outer extent is irrelevant. The beam
supports obey the same conclusion by their real centers and
exponentially thin tube margins.

Moreover $I(t_T)=T$ and $I'\le C_I$ show that the selected entry
points with $v<(1-\nu_0)T$ have
\[
 t_e\le t_T-\frac{\nu_0}{C_I}T+O(1).
\]
Choose $\alpha'<\nu_0/(2C_I)$ and make the first-hit allowance
smaller still. Ordinary local wave-map uniqueness, from exactly
equal metric and gauge Cauchy jets, proves exact equality on that
earlier entry interval. On $s=0$, the interval of
\eqref{pk2:domain} is $-u_2\le v\le T+c_2$; its negative part
is fixed, and its middle and unselected-end portions are protected
by finite-time domain of dependence. The deeper inner face is
future outflow, so no past influence can enter from the deep block.
There is consequently no additional dependence on an unexamined
route through the other end.

For the error estimate use the actual-principal ordinary scalar
stress energies on component errors, with lab mixed commutations
and a small fixed positive mass. The future slice normal is a
uniformly timelike multiplier, and both artificial boundaries
have strict outflow fluxes. The first- and second-order lab
background derivatives that determine the top linear rate are
globally bounded and tend uniformly to zero far out. Indeed the
linear tensor-wave lower terms are connections, their derivatives,
and curvature, while leading principal commutations use first
coefficient derivatives. Fix the mass small first, then choose
$R$ large: the square-root top rate on the far support is
arbitrarily small. On its complement that rate is bounded
independently of $R$ and $T$. Higher background hits multiply
lower commuted error norms, with subexponential coefficients.

The non-small-rate part of the support has time length at most
$\alpha'T+O(1)$ by \eqref{pk2:late-first-hit}. The lower-jet
triangular estimate used in the deep comparison costs just an
$e^{o(T)}$ factor. Retaining the exact principal coefficients,
quasilinear remainders have an exponentially small packet or
bootstrap factor against each remaining error placement, plus
the source in \eqref{pk2:beam-estimates}. Ordinary Sobolev product
estimates on the uniform annuli and inner charts apply without
an inverse lapse loss. Initial mixed jets follow from the equation
on uniformly spacelike slices. Taking the far rate and
$\alpha'$ sufficiently small therefore gives
\eqref{pk2:exterior-error}.

The parameter inequalities below strictly improve the exponentially
small bootstrap. Ordinary continuation preserves the two boundary
margins; exact wave-constraint propagation yields vacuum. Transform
the resulting entry jets to the fixed background wedge coordinates.
The transition and trace costs, with fixed buffers, are only
$e^{o(T)}$. The wave-map condition with target $g$ is tensorial and
exact. Beam matching adds the linear inverse-accuracy bound from
Proposition~\ref{pk2:beam-proposition}. The choices below make
the total smaller than \eqref{pk2:entry-accuracy}; the already
proved protected-region statement supplies its exact early part.
All estimates allow fixed small enlargements for the final gluing.
\end{proof}

\subsection{The order of the parameter choices}
\label{pk2:choices}

We verify simultaneously the comparison margins and approximation in
the prescribed data topology. First fix the finite data order $m$,
the correction order $s\ge m+3$, the interior and exterior comparison
orders, and every Sobolev, normal-jet, and trace buffer. Fix $N_b$
large enough for these requirements and then fix $P_*,P_0$ for the
quadratic residuals. None of these choices depends on formal accuracy.

Choose $\delta_2,\delta_3>0$ small relative to $\kappa$. The deep
mass, bulk threshold, and $\nu_0$ supply the rate in
\eqref{pk2:deep-error}; the far radius and localization parameters
$\alpha,\alpha'$ supply \eqref{pk2:exterior-error} and the exact
early agreement. The far radius may be enlarged later for
\eqref{pk2:arrival-shell}: the compact top rate is uniformly
bounded independently of it. These choices use fixed base bounds
and qualitative tails, with thresholds allowed to depend on the
smooth background datum.

Next take $\zeta>0$ small enough for every positive loss in the
chosen base norms, retaining more than $\kappa T/2$ of exponential
absorption for the perturbative deep products, and satisfying
\begin{align}
 \delta_3+(P_0-4)\zeta&<\kappa/2,&
 \delta_2+2\zeta&<\kappa/2,
 \label{pk2:positive-choices}\\
 \delta_2+(P_*-2)\zeta&<\kappa,&
 (m+\tfrac12)\zeta&<\kappa.
 \notag
\end{align}
Choose $0<b_p<\zeta/2$. Only now choose $M_1$ so large that
\begin{equation}
 (M_1+2)\zeta>\tfrac32\kappa+\delta_3,
 \qquad M_1\zeta>\kappa+\delta_2,
 \label{pk2:inverse-choices}
\end{equation}
and take sufficiently large finite transport, phase-Taylor, and
matching orders to achieve this accuracy. Finally take $T$ large
depending on all choices. Every subexponential factor can be
assigned a share of the strict margins.

By \eqref{pk2:packet-amplitude} and \eqref{pk2:blueshift},
$A_T=\exp((-2\kappa-2\zeta)T+o(T))$ with two-sided size.
The linear and quadratic exponents in the exterior error are
\[
 -2\kappa-(M_1+2)\zeta+\delta_3,
 \qquad -4\kappa+(P_0-4)\zeta+\delta_3.
\]
Both are strictly below $-7\kappa/2$ by
\eqref{pk2:positive-choices}--\eqref{pk2:inverse-choices}.
After division by $A_T$, the three terms of the deep error have
exponent bounds
\[
 -\tfrac12\kappa+\delta_2+2\zeta,
 \qquad \kappa+\delta_2-M_1\zeta,
 \qquad -\kappa+\delta_2+(P_*-2)\zeta,
\]
all strictly negative. This proves the entry condition and both
comparison improvements, and makes the observation errors
$o(Y_T^{-2})$. The direct beam contribution to the data seminorm
has exponent
\[
 -2\kappa+(m+\tfrac12)\zeta<0;
\]
the two correction contributions in
\eqref{pk2:data-approximation} also tend to zero. Thus
$p_m(d_T-d)\to0$ for this prescribed finite $m$. Smooth membership
of every $d_T$ was already obtained from the correction theorem;
no single frequency exponent has to give convergence at all orders.

\subsection{Attachment to the whole corrected bridge}
\label{pk2:attachment}

The finite exterior and deep solutions have now been constructed
with exact constraints, matching jets, and strict margins. It remains
to identify their observation with a point of the full MGHD of the
whole corrected datum. We attach these finite solutions to a lower
development containing every radius of the corrected bridge, then
use the outflow cuts to prove that the bridge remains Cauchy.

\begin{proposition}[Full-bridge Cauchy attachment]
\label{pk2:attachment-proposition}
For sufficiently large finite $T$, the exact data $d_T$ have a smooth
globally hyperbolic vacuum development containing the entire initial
bridge and the finite observation from
Lemma~\ref{pk2:observation}. Its adaptive backward geodesic reaches
$\Sigma$ with a seed $w_T\to w$. The development embeds openly
in the full MGHD of $d_T$, retaining the entire constraint-correction
tail on $\Sigma$.
\end{proposition}
\begin{proof}
The corrected data are complete and converge in $p_{10}$. They
agree with $d$ off the selected end and inside the protected ball.
The center satisfies $|z_*|\le s_*t_T+O(1)$, with
$s_*<1-2\alpha$, so that this ball contains the fixed-origin ball
of radius
\[
 (1-2\alpha-s_*)t_T-O(1)\longrightarrow\infty.
\]
Every fixed compact causal shadow used in the early bridge
construction is therefore unchanged for large $T$.

Construct the full finite lower bridge development for $d_T$,
including all initial spatial radii and a past collar, by the
finite-range Cauchy argument already proved. The bounded central
joins, the unselected end, and the bent compact part of the
selected $t_e=0$ slice agree with the original finite solution
by domain of dependence. Near infinity the selected horizontal
slice is the original initial surface carrying the actual corrected
data, including their tail.

Attach the selected ordinary wave-map exterior comparison through
an enlarged band about the fixed time $t_a$, only inside
\eqref{pk2:outer-cutoff} and before its terminal time. Use the
finite wave-map uniqueness and straightening from the comparison
development construction, with the same new data on both sides.
Fixed-time bounds and strict deeper outflow and spatial margins
give the identification, uniformly through the finite outer cutoff.
The straightening agrees on the used band and is cut off inside
its margin; on bounded central portions it is the original
identification because those solutions remain unchanged until
much later.

To see why the entire initial bridge is retained, put
$f(t_e,y)=|y|-R_T(t_e)$ and let $t_{\rm end}$ denote the comparison
terminal time. Locally the selected-end union is
\begin{equation}
 \{t_e<t_{\rm end}\}\cap
    \bigl(\{t_e<t_a\}\cup\{f<0\}\bigr),
 \label{pk2:partial-union}
\end{equation}
with the lower $\mathcal S$ cut imposed where active. The lower
set $t_e<t_a$ includes every spatial radius; only the added future
block is truncated. On future causal curves,
\[
 \frac{df}{dt_e}
   =\frac{d|y|}{dt_e}+V_{\rm out}
   \ge V_{\rm out}-V_{\rm max}>0.
\]
Thus the strict sublevels, their union, and their intersection
in \eqref{pk2:partial-union} are preserved to the past. A past
curve in the added block retains $f<0$ until it crosses $t_a$,
then enters the unrestricted lower block. It cannot end at the
omitted corner $t_e=t_a$, $f=0$. Unmatched future faces are omitted;
the original bridge is not truncated or replaced by the old datum
outside the computing cutoff.

The used portion of $\mathcal S$ and a matching collar lie within
the available margins. Its large positive endpoint has
$t_e\le t_T+O(1)$ and bounded co-radius, while the middle and
unselected portions are unchanged. The fixed background surface
is still spacelike for the new metric. Its exact comparison jets
and wave-map gauge agree with the deep Cauchy data. In the common
background collar, identify the deep vacuum solution across it,
retaining only the strict interior of \eqref{pk2:domain}.

More explicitly its active future cuts are
\[
 s<s_{\rm max},\qquad
 F_0=u+c_0s/T<u_2,\qquad
 F_1=v+c_0s/T<T+c_2,
 \qquad s_{\rm max}=(T+u_*)/2+c_0.
\]
If $\varepsilon_T$ is the exponentially small relative cone error,
then $du/ds,dv/ds\ge-C\varepsilon_T$ and hence
\[
 \frac{dF_i}{ds}\ge c_0/T-C\varepsilon_T>0.
\]
The signal point has $v=T$, $u=u_T<u_*$, with terminal margin
at least $c_0$ and strict side margins by the choice of $c_0$.
Past causal paths therefore reach $s=0$ strictly inside the
attached interval. For each finite $T$ all relevant coordinate
ranges are bounded and the metric is nondegenerate, so ordinary
continuation handles interior endpoints.

Omit every unused future face and edge from the lower pieces.
Their retained neighborhoods at intersections of cuts are finite
unions and intersections of the same past-preserved strict
sublevels; such corners cannot be past endpoints from the
interior. Use the identified collars to make the product gluing
of the comparison development, with no duplicate open
continuations. This yields a smooth Hausdorff vacuum manifold.
Every past-inextendible causal curve from the experiment reaches
the entry slice and then the complete corrected bridge. Lower
pieces have the finite Cauchy control already proved; bounded
speeds exclude escape to infinity in bounded time. The past
collar supplies the other direction of the Cauchy property.
Thus $\Sigma$ is Cauchy, and the maximal-development theorem
embeds this development into the full MGHD of $d_T$.

Finally the fixed earlier compact segment of the background
geodesic, through its transverse entry and back to $\Sigma$,
lies with a neighborhood in the unchanged region for large $T$.
The backward entry states from Lemma~\ref{pk2:observation}
continue through this fixed tube to seeds $w_T\to w$ by
ordinary geodesic flow dependence. This proves both the
attachment and the seed assertion.
\end{proof}

\subsection{Completion of the finite observation and genericity}
\label{pk2:conclusion}

\begin{proof}[Proof of Proposition~\ref{pk2:finite-signal}]
Fix $m$, $\eta$, and $J$. If a relatively open neighborhood
$\mathcal V$ of $d$ is also prescribed, choose a finite order
$m_{\mathcal V}$ and $\eta_{\mathcal V}>0$ such that the relative
ball $p_{m_{\mathcal V}}(d'-d)<\eta_{\mathcal V}$ lies in
$\mathcal V$. Replace the construction order by
$\widehat m=\max(m,m_{\mathcal V},10)$ and its tolerance by
$\widehat\eta=\min(\eta,\eta_{\mathcal V})$; without a prescribed
neighborhood use $\widehat m=m$ and $\widehat\eta=\eta$.
Make all choices in Subsection~\ref{pk2:choices} at this enlarged
finite order. The constrained beam correction
gives smooth complete exact vacuum data $d_T$ satisfying
\eqref{pk2:data-approximation}. The ordinary exterior comparison
supplies the exact entry jets required by the deep comparison.
Lemma~\ref{pk2:observation} gives the strict finite tidal
observation and convergent backward entry states.
Proposition~\ref{pk2:attachment-proposition} places the point,
its frame, and its geodesic back to $\Sigma$ in the full MGHD
of $d_T$, with $w_T\to w$. Convergence in $p_{\widehat m}$
attains both the original tolerance and the prescribed
neighborhood for sufficiently large finite $T$. The frequency
exponent may depend on this enlarged order; no convergence in
all seminorms for one fixed construction is asserted.
\end{proof}

\begin{proof}[Completion of the proof of Theorem~\ref{thm:main}]
Use the single base neighborhood constructed in the preceding
sections. Fix $O,N$, $d\in\mathcal B(O,N)$, a finite $m\ge10$,
and a tolerance $\eta>0$. The first-runout and corner-to-branch
arguments give a seed $w\in O$ with a noncorner first runout.
The closed test $A(O,N)$ supplies its bounded parallel curvature,
whether or not this new branch seed itself has a regular exit.

Choose a neighborhood of $(d,w)$ on which the bounded-seed tidal
comparison has a common constant $C_1$:
\[
 |\operatorname{Riem}(U',X',U',X')|
   \le C_1Q_{d'}(w',\tau)
 \quad\text{for }g_{d'}(X',X')=1,\quad g_{d'}(X',U')=0.
\]
Indeed parallel transport back to the initial frame gives the
explicit factor $(1+2h'(w',w'))^2$, uniformly bounded near $(d,w)$.
Choose $J>C_1N$ before constructing the late experiment. Apply
Proposition~\ref{pk2:finite-signal} with tolerance smaller than
\[
 \min\left\{\eta,
    \epsilon-p_{10}\bigl(d-(h_*,K_*)\bigr)\right\}.
\]
The second quantity is positive. For large finite $T$, one has
$w_T\in O$, $d_T\in\mathcal U_\epsilon$, $p_m(d_T-d)<\eta$, and, at
the finite point in its full MGHD,
\[
 Q_{d_T}(w_T,\tau_T)
 \ge C_1^{-1}
   |\operatorname{Riem}(g_{d_T})(U_T,X_T,U_T,X_T)|>N.
\]
This is exactly the destruction of the closed test required in
Section~\ref{sec:category}. Each relative closure of
$\mathcal B(O,N)$ is therefore nowhere dense. Their countable
union covers the future-extendible data, proving the stated
meagreness and the open-dense-intersection assertion.

We recall why the same $\epsilon$ serves all these tests. The
physical profile range and depth, computing depths, ray and
horizon margins are fixed first at the given mass and spin.
Next come the cylinder gauges, measuring radii, compact/far
absorption scales, fixed overlap times and regions, and the
base-order continuation thresholds for the exterior, entry,
and double-null constructions. All resulting base input
conditions on the two ends, bent slices, and finite joins follow
from one sufficiently small $p_{10}$ threshold. Higher smooth
estimates for a particular $d$ have datum- and order-dependent
constants, but arbitrarily small exponential rates obtained by
later-time or larger-radius choices, without higher-order data
smallness. The seed family, $N$, the finite $m$, branch location,
comparison orders, frequency, formal accuracy, and late time
affect only this finite experiment. They do not alter the base
neighborhood. All spin gaps and constants need be positive only
for the fixed $0<\mathfrak a<M$, as required.
\end{proof}
